%% SPDX-License-Identifier: CC-BY-4.0
%% Copyright (C) 2026  Alin-Adrian Anton <alin.anton@upt.ro>
%%
%% Shared with the sibling diode projects.  The canonical copy lives in
%% log-diode/docs/tutorial/ -- fix it there and re-copy, so opc-diode,
%% gnat-lt-pro, wal-postgresqldb and wal-mysqldb stay in step.
%% common.tex -- shared preamble for the diode tutorial documents.
%% Input from the preamble of both diode-quickstart.tex and
%% data-diode-handbook.tex.  Everything here must work with a stock
%% TeX Live / texlive-collection-latexrecommended install; no external
%% images, all figures are TikZ.

\usepackage[T1]{fontenc}
\usepackage{lmodern}
\usepackage{microtype}
\usepackage[a4paper,margin=25mm,headheight=15pt]{geometry}
\usepackage{amsmath,amssymb}
\usepackage{xcolor}
\usepackage{graphicx}      % \reflectbox (\copyleft), \resizebox for wide figures
\usepackage{booktabs,tabularx,longtable,array}
\usepackage{enumitem}
\usepackage{caption}
\usepackage{fancyhdr}
\usepackage{titlesec}
\usepackage{listings}
\usepackage{tikz}
\usetikzlibrary{positioning,arrows.meta,shapes.geometric,calc,fit,backgrounds}
\usepackage{tcolorbox}
\tcbuselibrary{skins,breakable}
\usepackage[hidelinks,pdfusetitle]{hyperref}

%% ---------------------------------------------------------------- colours
\definecolor{ldblue}  {HTML}{1F4E79}
\definecolor{ldred}   {HTML}{A61B1B}
\definecolor{ldgreen} {HTML}{1E6F3C}
\definecolor{ldamber} {HTML}{9A6700}
\definecolor{ldgray}  {HTML}{5A5A5A}
\definecolor{codebg}  {HTML}{F6F6F4}

%% ---------------------------------------------------------------- headers
\pagestyle{fancy}
\fancyhf{}
\fancyhead[L]{\footnotesize\textsc{\doctitleshort}}
\fancyhead[R]{\footnotesize\thepage}
\fancyfoot[C]{\footnotesize \projectname{} diode tutorial --- CC BY 4.0}
\renewcommand{\headrulewidth}{0.4pt}
\renewcommand{\footrulewidth}{0pt}

%% ---------------------------------------------------------------- listings
\lstset{
  basicstyle=\ttfamily\footnotesize,
  backgroundcolor=\color{codebg},
  frame=single,
  rulecolor=\color{ldgray!40},
  framesep=4pt,
  xleftmargin=6pt,
  xrightmargin=2pt,
  breaklines=true,
  breakatwhitespace=false,
  postbreak=\mbox{\textcolor{ldgray}{$\hookrightarrow$}\space},
  showstringspaces=false,
  keepspaces=true,
  columns=fullflexible,
  commentstyle=\color{ldgreen}\itshape,
  keywordstyle=\color{ldblue}\bfseries,
  stringstyle=\color{ldred},
  captionpos=b,
  aboveskip=8pt,
  belowskip=6pt,
  tabsize=2,
}

\lstdefinestyle{shell}{
  language=bash,
  morekeywords={ip,nft,iptables,ip6tables,pfctl,sysctl,systemctl,tcpdump,arp,
                ethtool,openssl,hping3,arping,nc,socat},
}
\lstdefinestyle{conf}{
  language=,
  morecomment=[l]{\#},
}
\lstdefinestyle{nft}{
  language=,
  morecomment=[l]{\#},
  morekeywords={table,chain,rule,type,hook,priority,policy,iif,oif,iifname,
                oifname,udp,dport,daddr,saddr,accept,drop,filter,ingress,
                input,output,forward,inet,netdev,ip,ip6,meta,counter,limit,log},
  keywordstyle=\color{ldblue}\bfseries,
}
\lstdefinestyle{pf}{
  language=,
  morecomment=[l]{\#},
  morekeywords={set,block,pass,in,out,quick,on,proto,from,to,port,all,
                no-state,state,antispoof,skip,block-policy,scrub,keep},
  keywordstyle=\color{ldblue}\bfseries,
}
\lstdefinestyle{ada}{
  language=Ada,
  morekeywords={SPARK_Mode,Pre,Post,Loop_Invariant},
}
\lstdefinestyle{c}{
  language=C,
  basicstyle=\ttfamily\scriptsize,
  morekeywords={ssize_t,size_t,sig_atomic_t,uint8_t,IFF_TAP,IFF_NO_PI,
                TUNSETIFF,POLLIN,O_RDWR,O_CLOEXEC,EINTR,EAGAIN,IFNAMSIZ},
  numbers=left,
  numberstyle=\ttfamily\tiny\color{ldgray!70},
  numbersep=7pt,
}

%% ------------------------------------------------------------- admonitions
\newtcolorbox{warnbox}[1][Warning]{%
  breakable, enhanced, colback=ldred!4, colframe=ldred,
  boxrule=0.9pt, left=6pt, right=6pt, top=5pt, bottom=5pt,
  fonttitle=\bfseries\sffamily\small, title={#1}}

\newtcolorbox{labbox}[1][Laboratory use only]{%
  breakable, enhanced, colback=ldamber!6, colframe=ldamber,
  boxrule=0.9pt, left=6pt, right=6pt, top=5pt, bottom=5pt,
  fonttitle=\bfseries\sffamily\small, title={#1}}

\newtcolorbox{notebox}[1][Note]{%
  breakable, enhanced, colback=ldblue!4, colframe=ldblue,
  boxrule=0.7pt, left=6pt, right=6pt, top=5pt, bottom=5pt,
  fonttitle=\bfseries\sffamily\small, title={#1}}

\newtcolorbox{tipbox}[1][In practice]{%
  breakable, enhanced, colback=ldgreen!4, colframe=ldgreen,
  boxrule=0.7pt, left=6pt, right=6pt, top=5pt, bottom=5pt,
  fonttitle=\bfseries\sffamily\small, title={#1}}

%% ------------------------------------------------------------------ macros
%% Copyleft (U+1F12F).  No Latin Modern / T1 glyph exists and neither
%% cclicenses nor ccicons is in a base TeX Live, so mirror the copyright
%% sign -- which is exactly what the symbol is.  \mathord keeps the spacing
%% sane when it lands next to text.
\newcommand{\copyleft}{\reflectbox{\copyright}}

\newcommand{\cmd}[1]{\texttt{#1}}
\newcommand{\file}[1]{\texttt{#1}}
\newcommand{\prog}[1]{\textsf{#1}}
%% Renewed per project, immediately after %% SPDX-License-Identifier: CC-BY-4.0
%% Copyright (C) 2026  Alin-Adrian Anton <alin.anton@upt.ro>
%%
%% Shared with the sibling diode projects.  The canonical copy lives in
%% log-diode/docs/tutorial/ -- fix it there and re-copy, so opc-diode,
%% gnat-lt-pro, wal-postgresqldb and wal-mysqldb stay in step.
%% common.tex -- shared preamble for the diode tutorial documents.
%% Input from the preamble of both diode-quickstart.tex and
%% data-diode-handbook.tex.  Everything here must work with a stock
%% TeX Live / texlive-collection-latexrecommended install; no external
%% images, all figures are TikZ.

\usepackage[T1]{fontenc}
\usepackage{lmodern}
\usepackage{microtype}
\usepackage[a4paper,margin=25mm,headheight=15pt]{geometry}
\usepackage{amsmath,amssymb}
\usepackage{xcolor}
\usepackage{graphicx}      % \reflectbox (\copyleft), \resizebox for wide figures
\usepackage{booktabs,tabularx,longtable,array}
\usepackage{enumitem}
\usepackage{caption}
\usepackage{fancyhdr}
\usepackage{titlesec}
\usepackage{listings}
\usepackage{tikz}
\usetikzlibrary{positioning,arrows.meta,shapes.geometric,calc,fit,backgrounds}
\usepackage{tcolorbox}
\tcbuselibrary{skins,breakable}
\usepackage[hidelinks,pdfusetitle]{hyperref}

%% ---------------------------------------------------------------- colours
\definecolor{ldblue}  {HTML}{1F4E79}
\definecolor{ldred}   {HTML}{A61B1B}
\definecolor{ldgreen} {HTML}{1E6F3C}
\definecolor{ldamber} {HTML}{9A6700}
\definecolor{ldgray}  {HTML}{5A5A5A}
\definecolor{codebg}  {HTML}{F6F6F4}

%% ---------------------------------------------------------------- headers
\pagestyle{fancy}
\fancyhf{}
\fancyhead[L]{\footnotesize\textsc{\doctitleshort}}
\fancyhead[R]{\footnotesize\thepage}
\fancyfoot[C]{\footnotesize \projectname{} diode tutorial --- CC BY 4.0}
\renewcommand{\headrulewidth}{0.4pt}
\renewcommand{\footrulewidth}{0pt}

%% ---------------------------------------------------------------- listings
\lstset{
  basicstyle=\ttfamily\footnotesize,
  backgroundcolor=\color{codebg},
  frame=single,
  rulecolor=\color{ldgray!40},
  framesep=4pt,
  xleftmargin=6pt,
  xrightmargin=2pt,
  breaklines=true,
  breakatwhitespace=false,
  postbreak=\mbox{\textcolor{ldgray}{$\hookrightarrow$}\space},
  showstringspaces=false,
  keepspaces=true,
  columns=fullflexible,
  commentstyle=\color{ldgreen}\itshape,
  keywordstyle=\color{ldblue}\bfseries,
  stringstyle=\color{ldred},
  captionpos=b,
  aboveskip=8pt,
  belowskip=6pt,
  tabsize=2,
}

\lstdefinestyle{shell}{
  language=bash,
  morekeywords={ip,nft,iptables,ip6tables,pfctl,sysctl,systemctl,tcpdump,arp,
                ethtool,openssl,hping3,arping,nc,socat},
}
\lstdefinestyle{conf}{
  language=,
  morecomment=[l]{\#},
}
\lstdefinestyle{nft}{
  language=,
  morecomment=[l]{\#},
  morekeywords={table,chain,rule,type,hook,priority,policy,iif,oif,iifname,
                oifname,udp,dport,daddr,saddr,accept,drop,filter,ingress,
                input,output,forward,inet,netdev,ip,ip6,meta,counter,limit,log},
  keywordstyle=\color{ldblue}\bfseries,
}
\lstdefinestyle{pf}{
  language=,
  morecomment=[l]{\#},
  morekeywords={set,block,pass,in,out,quick,on,proto,from,to,port,all,
                no-state,state,antispoof,skip,block-policy,scrub,keep},
  keywordstyle=\color{ldblue}\bfseries,
}
\lstdefinestyle{ada}{
  language=Ada,
  morekeywords={SPARK_Mode,Pre,Post,Loop_Invariant},
}
\lstdefinestyle{c}{
  language=C,
  basicstyle=\ttfamily\scriptsize,
  morekeywords={ssize_t,size_t,sig_atomic_t,uint8_t,IFF_TAP,IFF_NO_PI,
                TUNSETIFF,POLLIN,O_RDWR,O_CLOEXEC,EINTR,EAGAIN,IFNAMSIZ},
  numbers=left,
  numberstyle=\ttfamily\tiny\color{ldgray!70},
  numbersep=7pt,
}

%% ------------------------------------------------------------- admonitions
\newtcolorbox{warnbox}[1][Warning]{%
  breakable, enhanced, colback=ldred!4, colframe=ldred,
  boxrule=0.9pt, left=6pt, right=6pt, top=5pt, bottom=5pt,
  fonttitle=\bfseries\sffamily\small, title={#1}}

\newtcolorbox{labbox}[1][Laboratory use only]{%
  breakable, enhanced, colback=ldamber!6, colframe=ldamber,
  boxrule=0.9pt, left=6pt, right=6pt, top=5pt, bottom=5pt,
  fonttitle=\bfseries\sffamily\small, title={#1}}

\newtcolorbox{notebox}[1][Note]{%
  breakable, enhanced, colback=ldblue!4, colframe=ldblue,
  boxrule=0.7pt, left=6pt, right=6pt, top=5pt, bottom=5pt,
  fonttitle=\bfseries\sffamily\small, title={#1}}

\newtcolorbox{tipbox}[1][In practice]{%
  breakable, enhanced, colback=ldgreen!4, colframe=ldgreen,
  boxrule=0.7pt, left=6pt, right=6pt, top=5pt, bottom=5pt,
  fonttitle=\bfseries\sffamily\small, title={#1}}

%% ------------------------------------------------------------------ macros
%% Copyleft (U+1F12F).  No Latin Modern / T1 glyph exists and neither
%% cclicenses nor ccicons is in a base TeX Live, so mirror the copyright
%% sign -- which is exactly what the symbol is.  \mathord keeps the spacing
%% sane when it lands next to text.
\newcommand{\copyleft}{\reflectbox{\copyright}}

\newcommand{\cmd}[1]{\texttt{#1}}
\newcommand{\file}[1]{\texttt{#1}}
\newcommand{\prog}[1]{\textsf{#1}}
%% Renewed per project, immediately after %% SPDX-License-Identifier: CC-BY-4.0
%% Copyright (C) 2026  Alin-Adrian Anton <alin.anton@upt.ro>
%%
%% Shared with the sibling diode projects.  The canonical copy lives in
%% log-diode/docs/tutorial/ -- fix it there and re-copy, so opc-diode,
%% gnat-lt-pro, wal-postgresqldb and wal-mysqldb stay in step.
%% common.tex -- shared preamble for the diode tutorial documents.
%% Input from the preamble of both diode-quickstart.tex and
%% data-diode-handbook.tex.  Everything here must work with a stock
%% TeX Live / texlive-collection-latexrecommended install; no external
%% images, all figures are TikZ.

\usepackage[T1]{fontenc}
\usepackage{lmodern}
\usepackage{microtype}
\usepackage[a4paper,margin=25mm,headheight=15pt]{geometry}
\usepackage{amsmath,amssymb}
\usepackage{xcolor}
\usepackage{graphicx}      % \reflectbox (\copyleft), \resizebox for wide figures
\usepackage{booktabs,tabularx,longtable,array}
\usepackage{enumitem}
\usepackage{caption}
\usepackage{fancyhdr}
\usepackage{titlesec}
\usepackage{listings}
\usepackage{tikz}
\usetikzlibrary{positioning,arrows.meta,shapes.geometric,calc,fit,backgrounds}
\usepackage{tcolorbox}
\tcbuselibrary{skins,breakable}
\usepackage[hidelinks,pdfusetitle]{hyperref}

%% ---------------------------------------------------------------- colours
\definecolor{ldblue}  {HTML}{1F4E79}
\definecolor{ldred}   {HTML}{A61B1B}
\definecolor{ldgreen} {HTML}{1E6F3C}
\definecolor{ldamber} {HTML}{9A6700}
\definecolor{ldgray}  {HTML}{5A5A5A}
\definecolor{codebg}  {HTML}{F6F6F4}

%% ---------------------------------------------------------------- headers
\pagestyle{fancy}
\fancyhf{}
\fancyhead[L]{\footnotesize\textsc{\doctitleshort}}
\fancyhead[R]{\footnotesize\thepage}
\fancyfoot[C]{\footnotesize \projectname{} diode tutorial --- CC BY 4.0}
\renewcommand{\headrulewidth}{0.4pt}
\renewcommand{\footrulewidth}{0pt}

%% ---------------------------------------------------------------- listings
\lstset{
  basicstyle=\ttfamily\footnotesize,
  backgroundcolor=\color{codebg},
  frame=single,
  rulecolor=\color{ldgray!40},
  framesep=4pt,
  xleftmargin=6pt,
  xrightmargin=2pt,
  breaklines=true,
  breakatwhitespace=false,
  postbreak=\mbox{\textcolor{ldgray}{$\hookrightarrow$}\space},
  showstringspaces=false,
  keepspaces=true,
  columns=fullflexible,
  commentstyle=\color{ldgreen}\itshape,
  keywordstyle=\color{ldblue}\bfseries,
  stringstyle=\color{ldred},
  captionpos=b,
  aboveskip=8pt,
  belowskip=6pt,
  tabsize=2,
}

\lstdefinestyle{shell}{
  language=bash,
  morekeywords={ip,nft,iptables,ip6tables,pfctl,sysctl,systemctl,tcpdump,arp,
                ethtool,openssl,hping3,arping,nc,socat},
}
\lstdefinestyle{conf}{
  language=,
  morecomment=[l]{\#},
}
\lstdefinestyle{nft}{
  language=,
  morecomment=[l]{\#},
  morekeywords={table,chain,rule,type,hook,priority,policy,iif,oif,iifname,
                oifname,udp,dport,daddr,saddr,accept,drop,filter,ingress,
                input,output,forward,inet,netdev,ip,ip6,meta,counter,limit,log},
  keywordstyle=\color{ldblue}\bfseries,
}
\lstdefinestyle{pf}{
  language=,
  morecomment=[l]{\#},
  morekeywords={set,block,pass,in,out,quick,on,proto,from,to,port,all,
                no-state,state,antispoof,skip,block-policy,scrub,keep},
  keywordstyle=\color{ldblue}\bfseries,
}
\lstdefinestyle{ada}{
  language=Ada,
  morekeywords={SPARK_Mode,Pre,Post,Loop_Invariant},
}
\lstdefinestyle{c}{
  language=C,
  basicstyle=\ttfamily\scriptsize,
  morekeywords={ssize_t,size_t,sig_atomic_t,uint8_t,IFF_TAP,IFF_NO_PI,
                TUNSETIFF,POLLIN,O_RDWR,O_CLOEXEC,EINTR,EAGAIN,IFNAMSIZ},
  numbers=left,
  numberstyle=\ttfamily\tiny\color{ldgray!70},
  numbersep=7pt,
}

%% ------------------------------------------------------------- admonitions
\newtcolorbox{warnbox}[1][Warning]{%
  breakable, enhanced, colback=ldred!4, colframe=ldred,
  boxrule=0.9pt, left=6pt, right=6pt, top=5pt, bottom=5pt,
  fonttitle=\bfseries\sffamily\small, title={#1}}

\newtcolorbox{labbox}[1][Laboratory use only]{%
  breakable, enhanced, colback=ldamber!6, colframe=ldamber,
  boxrule=0.9pt, left=6pt, right=6pt, top=5pt, bottom=5pt,
  fonttitle=\bfseries\sffamily\small, title={#1}}

\newtcolorbox{notebox}[1][Note]{%
  breakable, enhanced, colback=ldblue!4, colframe=ldblue,
  boxrule=0.7pt, left=6pt, right=6pt, top=5pt, bottom=5pt,
  fonttitle=\bfseries\sffamily\small, title={#1}}

\newtcolorbox{tipbox}[1][In practice]{%
  breakable, enhanced, colback=ldgreen!4, colframe=ldgreen,
  boxrule=0.7pt, left=6pt, right=6pt, top=5pt, bottom=5pt,
  fonttitle=\bfseries\sffamily\small, title={#1}}

%% ------------------------------------------------------------------ macros
%% Copyleft (U+1F12F).  No Latin Modern / T1 glyph exists and neither
%% cclicenses nor ccicons is in a base TeX Live, so mirror the copyright
%% sign -- which is exactly what the symbol is.  \mathord keeps the spacing
%% sane when it lands next to text.
\newcommand{\copyleft}{\reflectbox{\copyright}}

\newcommand{\cmd}[1]{\texttt{#1}}
\newcommand{\file}[1]{\texttt{#1}}
\newcommand{\prog}[1]{\textsf{#1}}
%% Renewed per project, immediately after %% SPDX-License-Identifier: CC-BY-4.0
%% Copyright (C) 2026  Alin-Adrian Anton <alin.anton@upt.ro>
%%
%% Shared with the sibling diode projects.  The canonical copy lives in
%% log-diode/docs/tutorial/ -- fix it there and re-copy, so opc-diode,
%% gnat-lt-pro, wal-postgresqldb and wal-mysqldb stay in step.
%% common.tex -- shared preamble for the diode tutorial documents.
%% Input from the preamble of both diode-quickstart.tex and
%% data-diode-handbook.tex.  Everything here must work with a stock
%% TeX Live / texlive-collection-latexrecommended install; no external
%% images, all figures are TikZ.

\usepackage[T1]{fontenc}
\usepackage{lmodern}
\usepackage{microtype}
\usepackage[a4paper,margin=25mm,headheight=15pt]{geometry}
\usepackage{amsmath,amssymb}
\usepackage{xcolor}
\usepackage{graphicx}      % \reflectbox (\copyleft), \resizebox for wide figures
\usepackage{booktabs,tabularx,longtable,array}
\usepackage{enumitem}
\usepackage{caption}
\usepackage{fancyhdr}
\usepackage{titlesec}
\usepackage{listings}
\usepackage{tikz}
\usetikzlibrary{positioning,arrows.meta,shapes.geometric,calc,fit,backgrounds}
\usepackage{tcolorbox}
\tcbuselibrary{skins,breakable}
\usepackage[hidelinks,pdfusetitle]{hyperref}

%% ---------------------------------------------------------------- colours
\definecolor{ldblue}  {HTML}{1F4E79}
\definecolor{ldred}   {HTML}{A61B1B}
\definecolor{ldgreen} {HTML}{1E6F3C}
\definecolor{ldamber} {HTML}{9A6700}
\definecolor{ldgray}  {HTML}{5A5A5A}
\definecolor{codebg}  {HTML}{F6F6F4}

%% ---------------------------------------------------------------- headers
\pagestyle{fancy}
\fancyhf{}
\fancyhead[L]{\footnotesize\textsc{\doctitleshort}}
\fancyhead[R]{\footnotesize\thepage}
\fancyfoot[C]{\footnotesize \projectname{} diode tutorial --- CC BY 4.0}
\renewcommand{\headrulewidth}{0.4pt}
\renewcommand{\footrulewidth}{0pt}

%% ---------------------------------------------------------------- listings
\lstset{
  basicstyle=\ttfamily\footnotesize,
  backgroundcolor=\color{codebg},
  frame=single,
  rulecolor=\color{ldgray!40},
  framesep=4pt,
  xleftmargin=6pt,
  xrightmargin=2pt,
  breaklines=true,
  breakatwhitespace=false,
  postbreak=\mbox{\textcolor{ldgray}{$\hookrightarrow$}\space},
  showstringspaces=false,
  keepspaces=true,
  columns=fullflexible,
  commentstyle=\color{ldgreen}\itshape,
  keywordstyle=\color{ldblue}\bfseries,
  stringstyle=\color{ldred},
  captionpos=b,
  aboveskip=8pt,
  belowskip=6pt,
  tabsize=2,
}

\lstdefinestyle{shell}{
  language=bash,
  morekeywords={ip,nft,iptables,ip6tables,pfctl,sysctl,systemctl,tcpdump,arp,
                ethtool,openssl,hping3,arping,nc,socat},
}
\lstdefinestyle{conf}{
  language=,
  morecomment=[l]{\#},
}
\lstdefinestyle{nft}{
  language=,
  morecomment=[l]{\#},
  morekeywords={table,chain,rule,type,hook,priority,policy,iif,oif,iifname,
                oifname,udp,dport,daddr,saddr,accept,drop,filter,ingress,
                input,output,forward,inet,netdev,ip,ip6,meta,counter,limit,log},
  keywordstyle=\color{ldblue}\bfseries,
}
\lstdefinestyle{pf}{
  language=,
  morecomment=[l]{\#},
  morekeywords={set,block,pass,in,out,quick,on,proto,from,to,port,all,
                no-state,state,antispoof,skip,block-policy,scrub,keep},
  keywordstyle=\color{ldblue}\bfseries,
}
\lstdefinestyle{ada}{
  language=Ada,
  morekeywords={SPARK_Mode,Pre,Post,Loop_Invariant},
}
\lstdefinestyle{c}{
  language=C,
  basicstyle=\ttfamily\scriptsize,
  morekeywords={ssize_t,size_t,sig_atomic_t,uint8_t,IFF_TAP,IFF_NO_PI,
                TUNSETIFF,POLLIN,O_RDWR,O_CLOEXEC,EINTR,EAGAIN,IFNAMSIZ},
  numbers=left,
  numberstyle=\ttfamily\tiny\color{ldgray!70},
  numbersep=7pt,
}

%% ------------------------------------------------------------- admonitions
\newtcolorbox{warnbox}[1][Warning]{%
  breakable, enhanced, colback=ldred!4, colframe=ldred,
  boxrule=0.9pt, left=6pt, right=6pt, top=5pt, bottom=5pt,
  fonttitle=\bfseries\sffamily\small, title={#1}}

\newtcolorbox{labbox}[1][Laboratory use only]{%
  breakable, enhanced, colback=ldamber!6, colframe=ldamber,
  boxrule=0.9pt, left=6pt, right=6pt, top=5pt, bottom=5pt,
  fonttitle=\bfseries\sffamily\small, title={#1}}

\newtcolorbox{notebox}[1][Note]{%
  breakable, enhanced, colback=ldblue!4, colframe=ldblue,
  boxrule=0.7pt, left=6pt, right=6pt, top=5pt, bottom=5pt,
  fonttitle=\bfseries\sffamily\small, title={#1}}

\newtcolorbox{tipbox}[1][In practice]{%
  breakable, enhanced, colback=ldgreen!4, colframe=ldgreen,
  boxrule=0.7pt, left=6pt, right=6pt, top=5pt, bottom=5pt,
  fonttitle=\bfseries\sffamily\small, title={#1}}

%% ------------------------------------------------------------------ macros
%% Copyleft (U+1F12F).  No Latin Modern / T1 glyph exists and neither
%% cclicenses nor ccicons is in a base TeX Live, so mirror the copyright
%% sign -- which is exactly what the symbol is.  \mathord keeps the spacing
%% sane when it lands next to text.
\newcommand{\copyleft}{\reflectbox{\copyright}}

\newcommand{\cmd}[1]{\texttt{#1}}
\newcommand{\file}[1]{\texttt{#1}}
\newcommand{\prog}[1]{\textsf{#1}}
%% Renewed per project, immediately after \input{common}.
\newcommand{\amp}{\prog{sender}}
\newcommand{\deamp}{\prog{receiver}}
%% The two ends of the link, named by ROLE rather than by security level.
%% A diode can be deployed in either orientation -- ingesting into a
%% protected domain, or exporting out of one -- so the mechanics here are
%% written in terms of which end sends and which receives. Each project's
%% "Orientation" section maps these onto its own security levels.
\newcommand{\srcside}{\textsc{source}}
\newcommand{\dstside}{\textsc{destination}}
\newcommand{\nm}[1]{#1\,nm}
\newcommand{\dbm}[1]{#1\,dBm}
\newcommand{\dbb}[1]{#1\,dB}

%% ------------------------------------------------------------ tikz styles
\tikzset{
  host/.style     ={draw, rounded corners=2pt, minimum width=27mm,
                    minimum height=12mm, align=center, font=\small,
                    fill=ldblue!8, draw=ldblue!70},
  mc/.style       ={draw, minimum width=25mm, minimum height=13mm,
                    align=center, font=\scriptsize, fill=gray!8,
                    draw=ldgray},
  netbox/.style   ={draw, dashed, rounded corners=3pt, draw=ldgray!70,
                    inner sep=5mm},
  splitter/.style ={draw, diamond, aspect=1.5, inner sep=1.5pt,
                    align=center, font=\scriptsize, fill=ldgreen!12,
                    draw=ldgreen},
  fiber/.style    ={-{Latex[length=2.4mm,width=1.8mm]}, thick,
                    draw=ldgreen!65!black},
  copper/.style   ={-, thick, draw=ldblue!70!black},
  dataflow/.style ={-{Latex[length=3.4mm,width=2.6mm]}, line width=1.5pt,
                    draw=ldblue},
  blocked/.style  ={-{Latex[length=2.4mm]}, thick, draw=ldred, dashed},
  lbl/.style      ={font=\scriptsize\sffamily, inner sep=1.5pt},
  lblg/.style     ={font=\scriptsize\sffamily, inner sep=1.5pt,
                    text=ldgreen!55!black},
  lblr/.style     ={font=\scriptsize\sffamily, inner sep=1.5pt, text=ldred},
}

%% A red "no return path" marker, placed on a path with node[midway]{\xmark}.
\newcommand{\xmark}{\textcolor{ldred}{\sffamily\bfseries\large $\times$}}

%% ------------------------------------------------------------ title block
\titleformat{\section}{\normalfont\Large\bfseries\color{ldblue}}{\thesection}{0.7em}{}
\titleformat{\subsection}{\normalfont\large\bfseries\color{ldblue!85!black}}{\thesubsection}{0.6em}{}
\titleformat{\subsubsection}{\normalfont\normalsize\bfseries}{\thesubsubsection}{0.6em}{}
\setlist[itemize]{itemsep=2pt, topsep=3pt, parsep=0pt}
\setlist[enumerate]{itemsep=2pt, topsep=3pt, parsep=0pt}
\captionsetup{font=small, labelfont={bf,color=ldblue}}
.
\newcommand{\amp}{\prog{sender}}
\newcommand{\deamp}{\prog{receiver}}
%% The two ends of the link, named by ROLE rather than by security level.
%% A diode can be deployed in either orientation -- ingesting into a
%% protected domain, or exporting out of one -- so the mechanics here are
%% written in terms of which end sends and which receives. Each project's
%% "Orientation" section maps these onto its own security levels.
\newcommand{\srcside}{\textsc{source}}
\newcommand{\dstside}{\textsc{destination}}
\newcommand{\nm}[1]{#1\,nm}
\newcommand{\dbm}[1]{#1\,dBm}
\newcommand{\dbb}[1]{#1\,dB}

%% ------------------------------------------------------------ tikz styles
\tikzset{
  host/.style     ={draw, rounded corners=2pt, minimum width=27mm,
                    minimum height=12mm, align=center, font=\small,
                    fill=ldblue!8, draw=ldblue!70},
  mc/.style       ={draw, minimum width=25mm, minimum height=13mm,
                    align=center, font=\scriptsize, fill=gray!8,
                    draw=ldgray},
  netbox/.style   ={draw, dashed, rounded corners=3pt, draw=ldgray!70,
                    inner sep=5mm},
  splitter/.style ={draw, diamond, aspect=1.5, inner sep=1.5pt,
                    align=center, font=\scriptsize, fill=ldgreen!12,
                    draw=ldgreen},
  fiber/.style    ={-{Latex[length=2.4mm,width=1.8mm]}, thick,
                    draw=ldgreen!65!black},
  copper/.style   ={-, thick, draw=ldblue!70!black},
  dataflow/.style ={-{Latex[length=3.4mm,width=2.6mm]}, line width=1.5pt,
                    draw=ldblue},
  blocked/.style  ={-{Latex[length=2.4mm]}, thick, draw=ldred, dashed},
  lbl/.style      ={font=\scriptsize\sffamily, inner sep=1.5pt},
  lblg/.style     ={font=\scriptsize\sffamily, inner sep=1.5pt,
                    text=ldgreen!55!black},
  lblr/.style     ={font=\scriptsize\sffamily, inner sep=1.5pt, text=ldred},
}

%% A red "no return path" marker, placed on a path with node[midway]{\xmark}.
\newcommand{\xmark}{\textcolor{ldred}{\sffamily\bfseries\large $\times$}}

%% ------------------------------------------------------------ title block
\titleformat{\section}{\normalfont\Large\bfseries\color{ldblue}}{\thesection}{0.7em}{}
\titleformat{\subsection}{\normalfont\large\bfseries\color{ldblue!85!black}}{\thesubsection}{0.6em}{}
\titleformat{\subsubsection}{\normalfont\normalsize\bfseries}{\thesubsubsection}{0.6em}{}
\setlist[itemize]{itemsep=2pt, topsep=3pt, parsep=0pt}
\setlist[enumerate]{itemsep=2pt, topsep=3pt, parsep=0pt}
\captionsetup{font=small, labelfont={bf,color=ldblue}}
.
\newcommand{\amp}{\prog{sender}}
\newcommand{\deamp}{\prog{receiver}}
%% The two ends of the link, named by ROLE rather than by security level.
%% A diode can be deployed in either orientation -- ingesting into a
%% protected domain, or exporting out of one -- so the mechanics here are
%% written in terms of which end sends and which receives. Each project's
%% "Orientation" section maps these onto its own security levels.
\newcommand{\srcside}{\textsc{source}}
\newcommand{\dstside}{\textsc{destination}}
\newcommand{\nm}[1]{#1\,nm}
\newcommand{\dbm}[1]{#1\,dBm}
\newcommand{\dbb}[1]{#1\,dB}

%% ------------------------------------------------------------ tikz styles
\tikzset{
  host/.style     ={draw, rounded corners=2pt, minimum width=27mm,
                    minimum height=12mm, align=center, font=\small,
                    fill=ldblue!8, draw=ldblue!70},
  mc/.style       ={draw, minimum width=25mm, minimum height=13mm,
                    align=center, font=\scriptsize, fill=gray!8,
                    draw=ldgray},
  netbox/.style   ={draw, dashed, rounded corners=3pt, draw=ldgray!70,
                    inner sep=5mm},
  splitter/.style ={draw, diamond, aspect=1.5, inner sep=1.5pt,
                    align=center, font=\scriptsize, fill=ldgreen!12,
                    draw=ldgreen},
  fiber/.style    ={-{Latex[length=2.4mm,width=1.8mm]}, thick,
                    draw=ldgreen!65!black},
  copper/.style   ={-, thick, draw=ldblue!70!black},
  dataflow/.style ={-{Latex[length=3.4mm,width=2.6mm]}, line width=1.5pt,
                    draw=ldblue},
  blocked/.style  ={-{Latex[length=2.4mm]}, thick, draw=ldred, dashed},
  lbl/.style      ={font=\scriptsize\sffamily, inner sep=1.5pt},
  lblg/.style     ={font=\scriptsize\sffamily, inner sep=1.5pt,
                    text=ldgreen!55!black},
  lblr/.style     ={font=\scriptsize\sffamily, inner sep=1.5pt, text=ldred},
}

%% A red "no return path" marker, placed on a path with node[midway]{\xmark}.
\newcommand{\xmark}{\textcolor{ldred}{\sffamily\bfseries\large $\times$}}

%% ------------------------------------------------------------ title block
\titleformat{\section}{\normalfont\Large\bfseries\color{ldblue}}{\thesection}{0.7em}{}
\titleformat{\subsection}{\normalfont\large\bfseries\color{ldblue!85!black}}{\thesubsection}{0.6em}{}
\titleformat{\subsubsection}{\normalfont\normalsize\bfseries}{\thesubsubsection}{0.6em}{}
\setlist[itemize]{itemsep=2pt, topsep=3pt, parsep=0pt}
\setlist[enumerate]{itemsep=2pt, topsep=3pt, parsep=0pt}
\captionsetup{font=small, labelfont={bf,color=ldblue}}
.
\newcommand{\amp}{\prog{sender}}
\newcommand{\deamp}{\prog{receiver}}
%% The two ends of the link, named by ROLE rather than by security level.
%% A diode can be deployed in either orientation -- ingesting into a
%% protected domain, or exporting out of one -- so the mechanics here are
%% written in terms of which end sends and which receives. Each project's
%% "Orientation" section maps these onto its own security levels.
\newcommand{\srcside}{\textsc{source}}
\newcommand{\dstside}{\textsc{destination}}
\newcommand{\nm}[1]{#1\,nm}
\newcommand{\dbm}[1]{#1\,dBm}
\newcommand{\dbb}[1]{#1\,dB}

%% ------------------------------------------------------------ tikz styles
\tikzset{
  host/.style     ={draw, rounded corners=2pt, minimum width=27mm,
                    minimum height=12mm, align=center, font=\small,
                    fill=ldblue!8, draw=ldblue!70},
  mc/.style       ={draw, minimum width=25mm, minimum height=13mm,
                    align=center, font=\scriptsize, fill=gray!8,
                    draw=ldgray},
  netbox/.style   ={draw, dashed, rounded corners=3pt, draw=ldgray!70,
                    inner sep=5mm},
  splitter/.style ={draw, diamond, aspect=1.5, inner sep=1.5pt,
                    align=center, font=\scriptsize, fill=ldgreen!12,
                    draw=ldgreen},
  fiber/.style    ={-{Latex[length=2.4mm,width=1.8mm]}, thick,
                    draw=ldgreen!65!black},
  copper/.style   ={-, thick, draw=ldblue!70!black},
  dataflow/.style ={-{Latex[length=3.4mm,width=2.6mm]}, line width=1.5pt,
                    draw=ldblue},
  blocked/.style  ={-{Latex[length=2.4mm]}, thick, draw=ldred, dashed},
  lbl/.style      ={font=\scriptsize\sffamily, inner sep=1.5pt},
  lblg/.style     ={font=\scriptsize\sffamily, inner sep=1.5pt,
                    text=ldgreen!55!black},
  lblr/.style     ={font=\scriptsize\sffamily, inner sep=1.5pt, text=ldred},
}

%% A red "no return path" marker, placed on a path with node[midway]{\xmark}.
\newcommand{\xmark}{\textcolor{ldred}{\sffamily\bfseries\large $\times$}}

%% ------------------------------------------------------------ title block
\titleformat{\section}{\normalfont\Large\bfseries\color{ldblue}}{\thesection}{0.7em}{}
\titleformat{\subsection}{\normalfont\large\bfseries\color{ldblue!85!black}}{\thesubsection}{0.6em}{}
\titleformat{\subsubsection}{\normalfont\normalsize\bfseries}{\thesubsubsection}{0.6em}{}
\setlist[itemize]{itemsep=2pt, topsep=3pt, parsep=0pt}
\setlist[enumerate]{itemsep=2pt, topsep=3pt, parsep=0pt}
\captionsetup{font=small, labelfont={bf,color=ldblue}}
